\documentclass[11pt]{article}
\usepackage[utf8]{inputenc}
\usepackage[T2A]{fontenc}
\usepackage[russian]{babel}
\usepackage{amsmath,amssymb}
\usepackage[margin=2.5cm]{geometry}
\usepackage{hyperref}

\title{Топологии взаимодействия между роботами}
\author{}
\date{11 сентября 2026}

\begin{document}
\maketitle

\section{Почему это не peer-to-peer по умолчанию}

Взаимодействие между роботами флота почти всегда идёт не напрямую, а через
посредника — по трём причинам сразу: ограниченная пропускная способность,
сетевые ограничения (NAT, разные домашние сети), и, что важнее всего для
privacy-анализа, модель угроз становится проще при одном координаторе, чем при
$N \times N$ связях между роботами.

\section{Наблюдаемые топологии}

\begin{description}
  \item[Звезда (star).] Все роботы взаимодействуют с одним сервером/агрегатором;
        друг друга не видят и не связываются напрямую. Доминирующий случай на
        практике; ровно эта топология используется как threat model по умолчанию
        в security-роадмапе проекта (агрегатор = honest-but-curious сервер).
  \item[Иерархическая звезда.] Роботы $\to$ edge-координатор (локальный, в том же
        доме/объекте) $\to$ облачный сервер. Снижает задержку и трафик, добавляет
        промежуточную точку доверия и потенциальной атаки (см.\ раздел про методы
        агрегации).
  \item[Децентрализованная / gossip FL.] Роботы обмениваются обновлениями попарно
        (соседи по графу связи), без центрального сервера вообще; каждый постепенно
        сходится к консенсусной модели \cite{topologysurvey2023,decentralizedsurvey2023}.
        Убирает единую точку отказа/наблюдения, но коммуникационно дороже звезды и
        сходится медленнее — обновления распространяются по графу постепенно, а
        синхронизация усложняется разной скоростью обработки и качеством связи узлов.
  \item[Gossip через random walk.] Модель \guillemotleft путешествует\guillemotright{}
        по графу клиентов вместо broadcast всем соседям \cite{randomwalk2026},
        снижая коммуникационную нагрузку ценой более медленной сходимости; выбор
        конкретных соседей для обмена можно дополнительно оптимизировать
        \cite{graphgossip2026}.
  \item[Policy merging (без взаимодействия во время обучения).] Роботы обучаются
        полностью автономно, политики сливаются постфактум с учётом перестановочной
        инвариантности контроллеров \cite{fleetmerge2023}.
\end{description}

\section{Приватность при меняющемся графе связи}

Работа \emph{Privacy-preserving Decentralized Federated Learning over Time-varying
Communication Graph} \cite{timevarying2022} напрямую релевантна флоту роботов: граф
связей между устройствами может меняться во времени (робот переехал, сменил
Wi-Fi-сеть) — гарантия приватности должна быть устойчива и к этой динамике, не
только к статичному графу.

Секьюрная decentralized FL через gossip и virtual voting \cite{gossipvoting2026}
совмещает gossip-топологию с verifiable consensus для устойчивости к
злонамеренным узлам без центрального сервера — альтернатива Shamir-based secure
aggregation в топологиях без единого координатора.

\section{Следствие для модели угроз проекта}

Топология напрямую определяет threat model. В звезде атакующий — это сервер
(gray-box, видит все обновления) или внешний наблюдатель с доступом к развёрнутой
политике (black-box, видит только действия). В gossip/decentralized режиме модель
угроз меняется качественно: скомпрометированный сосед-робот видит только соседние
обновления, а не весь флот, что одновременно хуже (больше потенциальных точек атаки
в сети) и лучше (ни один участник не видит агрегированную картину целиком).

Практическая рекомендация для проекта: зафиксировать звезду с промежуточным
edge-узлом как основную топологию (самая частая в найденных robot-специфичных
источниках 2026 года, самая простая для формализации threat model), а перенос
результата на decentralized gossip оставить отдельным, явно помеченным
расширением с собственной моделью угроз.

\section{Тезаурус}

\begin{description}
  \item[Звезда (star topology)] Топология, в которой все клиенты взаимодействуют
    только с одним центральным сервером и не связываются друг с другом напрямую.
  \item[Иерархическая звезда] Звезда с промежуточным уровнем: клиенты $\to$
    локальный edge-координатор $\to$ центральный сервер.
  \item[Decentralized / gossip FL] Топология без центрального сервера: клиенты
    обмениваются обновлениями попарно с соседями по графу связи, постепенно
    сходясь к общей модели.
  \item[Граф связи] Структура, описывающая, какие пары клиентов могут напрямую
    обмениваться сообщениями в децентрализованной схеме.
  \item[Random walk (агрегация)] Вариант gossip, в котором модель передаётся от
    клиента к клиенту последовательно (как при случайном блуждании по графу), а не
    рассылается сразу всем соседям.
  \item[Virtual voting] Механизм консенсуса поверх gossip-сети, позволяющий
    участникам согласовать корректный результат агрегации без центрального сервера
    и без взаимного доверия.
  \item[Time-varying communication graph] Граф связи, состав рёбер которого меняется
    со временем (устройство сменило сеть, переместилось) — приватность и сходимость
    должны быть устойчивы к этой динамике.
  \item[Policy merging] Слияние весов независимо обученных политик постфактум, без
    какого-либо online-взаимодействия роботов во время обучения.
\end{description}

\begin{thebibliography}{9}
\bibitem{topologysurvey2023} \emph{Topology-aware Federated Learning in Edge
  Computing: A Comprehensive Survey}. arXiv:2302.02573, 2023.
\bibitem{decentralizedsurvey2023} \emph{Decentralized Federated Learning: A Survey
  and Perspective}. arXiv:2306.01603, 2023.
\bibitem{randomwalk2026} \emph{Decentralized Federated Averaging via Random Walk}.
  arXiv:2508.21286, 2026.
\bibitem{graphgossip2026} \emph{Graph-based Gossiping for Communication Efficiency
  in Decentralized Federated Learning}. arXiv:2506.10607, 2026.
\bibitem{fleetmerge2023} \emph{Robot Fleet Learning via Policy Merging}.
  arXiv:2310.01362, 2023.
\bibitem{timevarying2022} \emph{Privacy-preserving Decentralized Federated Learning
  over Time-varying Communication Graph}. arXiv:2210.00325, 2022.
\bibitem{gossipvoting2026} \emph{Secure Decentralized Federated Learning via Gossip
  and Virtual Voting}. arXiv:2607.08651, 2026.
\end{thebibliography}

\end{document}
